\clearpage\SourceLeaf{12}
\begin{center}\bfseries [NOTIZEN ÜBER KUBISCHE UND BIQUADRATISCHE RESTE.]\end{center}
\begin{center}[1.]\par\smallskip Eine höchst elegante, durch Induktion gewonnene Beobachtung.\end{center}
\noindent\hspace*{\parindent}2 ist ein kubischer Rest oder Nichtrest der Primzahl $p$ von der Form $3n+1$, je nachdem, ob $p$ darstellbar ist durch die Form
\[xx+27yy\qquad\text{oder}\qquad4xx+2xy+7yy.\]
\par\noindent\hspace*{\parindent}3 ist ein Rest oder Nichtrest, je nachdem, ob $p$ darstellbar ist durch
\[\begin{aligned} &xx+243yy\quad\text{oder}\quad4xx+2xy+61yy\\
\text{oder}\quad&7xx+6xy+36yy\quad\text{oder}\quad9xx+6xy+28yy.\end{aligned}\]
\par\noindent\hspace*{\parindent}5 ist ein $\left.\begin{array}{l}\text{Rest}\\\text{Nichtrest}\end{array}\right\}$, wenn $p$ dargestellt wird
\begingroup\small
\[\text{durch }\left\{\begin{array}{l}(1,0,675),\ (25,0,27),\ (13,1,52),\ (4,1,169)\\
(7,2,97),\ (9,3,76),\ (19,3,36),\ (25,5,28),\ (25,10,31),\ (27,9,28).\end{array}\right.\]
\endgroup
\begin{center}[2.]\end{center}
\noindent\hspace*{\parindent}Ein Kriterium, mit dem man entscheidet, ob die Zahl $m$ ein kubischer Rest der Primzahl $p$ von der Form $3n+1$ ist.

Es sei
\[4p=aa+27bb\]
